\section{Experimental Protocol}
\label{sec:protocol}

\subsection{Component-Wise Evaluation Design}

The study uses two complementary benchmarks because no retained cohort contains
all required static features, dynamic sequences, active-learning identity, and
fidelity artifacts. MalBehavD-V1 answers RQ1. An identity-audited EMBER subset
answers RQ2 and RQ3. Table~\ref{tab:data} makes their semantics and roles
explicit. Results are never pooled across datasets.

\begin{table*}[t]
\centering
\caption{Dataset provenance, label semantics, observation unit, and evaluation role.}
\label{tab:data}
\scriptsize
\setlength{\tabcolsep}{3pt}
\begin{tabular}{p{2.5cm}p{2.2cm}p{2.2cm}p{3.2cm}p{2.2cm}p{2.9cm}}
\toprule
Dataset & Retained samples & Class counts & Inputs / observation unit & Split & Role \\
\midrule
MalBehavD-V1 & 1,922 usable of 1,923 with $\geq20$ events & 932 benign / 990 malware & 199 static metadata attributes and ordered API-event prefixes & 5 outer folds; seed 2026 & Fusion baselines and P/PI/PID ablation (RQ1) \\
EMBER-derived identity audit & 2,000 & 1,500 benign / 500 malware & 2,381 static PE features; unique SHA-256, row ID, source index & 5 random outer folds; 5 AL/model seeds & Active learning (RQ2) and TreeSHAP audit (RQ3) \\
MalBehavD-V1 horizon-40 subset & 1,186 & Generic benign/malware & Ordered API-event prefixes through event 40 & Development subset & Exploratory robust PID only \\
Timestamp feasibility pilot & 8 usable sequences from 20 selected samples & 2 benign / 6 malware usable & 4,697 API--timestamp pairs from Cuckoofork reports & No outer-fold evaluation & Parser/data-source feasibility only \\
\bottomrule
\end{tabular}
\end{table*}

Both labels mean generic benign/malware status. AVClass is available for 498
EMBER malware rows but is not used as ransomware ground truth. MalBehavD-V1 may
contain source differences between benign and malware collection. The experiment
does not infer family-held-out, temporal, or cross-organization performance.

\subsection{EMBER Identity, Splits, and Configuration}

The EMBER-derived table contains 2,000 unique SHA-256 identities, unique row and
source indices, and no exactly duplicated feature vectors. Samples appeared in
November--December 2018. The locked dataset SHA-256 is
\texttt{69f0ae680846edd95458d1e7aaabe36f8}\-
\texttt{866f4f22a6d11488b126ed026b28963}, and the manifest SHA-256 is
\texttt{a2ca16402653e6cc51fdd7a4857d007f}\-
\texttt{af86807b044d78654f180409c73565bb}. The subset was recovered from EMBER
material but does not reproduce the complete official benchmark protocol.

Five random outer folds create held-out test partitions. Within each outer
development partition, retained rows are separated into an acquisition pool and
validation resource. Test rows do not enter query selection, target definition,
threshold choice, or stopping. All policies in the same fold/seed receive the
same initial labeled set and remaining pool. Exact identity auditing reduces
row-level leakage, but random folds can still place related families, sources,
or time periods on both sides.

Each AL run begins with 200 stratified labels and performs ten rounds of 100
queries, reaching a maximum of 1,200 training labels. Five AL seeds are used per
outer fold. The retained LightGBM configuration uses 60 trees, learning rate
0.05, 15 leaves, maximum depth 4, and seed-specific deterministic training. Each
round retains pool membership, query components, selected row IDs, model
configuration, and validation/test predictions.

For fidelity, each fold fits five LightGBM models with seeds 41--45, 100 trees,
learning rate 0.1, feature fraction 0.8, bagging fraction 0.8, 15 leaves, and
maximum depth 4. The 25 fitted models are evaluated with mean and median
replacement, three deletion budgets, and 20 random masks per budget. Seed
comparisons stay within a fold; fold comparisons change the data partition.

\subsection{MalBehavD-V1 Pairing and Fusion Protocol}

The MalBehavD-V1 source contains 2,570 rows, balanced between 1,285 benign and
1,285 malware entries, with 2,543 unique SHA-256 values. Twenty-six SHA values
are duplicated without label conflict. Requiring at least 20 API events yields
1,923 candidates; 1,922 have usable static metadata and dynamic sequences,
providing 99.95\% metadata coverage. The final paired cohort contains 932 benign
and 990 malware samples. Its manifest SHA-256 is
\texttt{9540b88b0da2463f63b4e983f5ecc277}\-
\texttt{e3845b7badebbb79695e35f2c3b1a669}.

Static input comprises 199 structural metadata attributes. The allow-list
excludes verdict, detection count, threat name, tag, and filename fields that
could directly encode collection or label information. The dynamic input is an
ordered API sequence scored at prefixes. The collector obtains metadata/report
records by SHA-256 and does not download or execute PE objects. Event index is
the primary observation axis; it cannot support a wall-clock latency claim.

Five outer folds use seed 2026. Model, fusion parameters, fixed weight, stacking,
and decision thresholds are selected with development/validation resources and
then locked for each held-out evaluation. All fusion methods use the same paired
channel scores. Paired bootstrap intervals compare predictions on identical
rows. The threshold objective is validation based, so every reported FPR is
paired with TPR and is not treated as a common-TPR comparison unless explicitly
stated.

The original confirmatory run established channel complementarity and the
absence of a useful derivative effect under its protocol. Adaptive P and the
new fusion-only comparison were then developed on the same 1,922-sample cohort.
Those values are measured on outer-test rows but are classified as development
evidence because design choices followed inspection of the confirmatory result.
External or temporal validation is required before promoting them to an
independent confirmatory claim.

\subsection{Metrics and Paired Reporting}

Fusion reports ROC--AUC, AP, FPR, TPR, and Brier score. The main table compares
Average, Max-risk, Fixed, Logistic stacking, Adaptive P, Adaptive PI, and Full
PID. Primary paired effects are Adaptive P minus Fixed and Full PID minus
Logistic stacking. An additional P versus PID comparison tests incremental I/D
value.

Active learning reports ROC-AULC, AP-AULC, training labels to the fold-specific
target, restricted-FPR summaries, and scoped computation. The paired comparison
unit is outer fold after averaging its seed trajectories; 25 fold--seed runs are
not treated as 25 independent datasets. Query timing includes pool prediction,
TreeSHAP where applicable, additivity verification, normalization, and ranking.
It excludes human annotation and endpoint feature extraction.

Fidelity reports AUC decrease under top and random deletion, the absolute
top-minus-random effect, Spearman rank correlation, and top-$k$ Jaccard
agreement. The 20 random masks describe control variation conditional on fitted
models; they are not independent deployment datasets. EC, CF, and additivity are
secondary diagnostics and are not combined into a composite score.

\subsection{Confirmatory versus Exploratory Status}

Table~\ref{tab:evidence-tier} locks the evidence hierarchy before interpretation.
This prevents a favorable exploratory value from silently replacing a less
favorable confirmatory result.

\begin{table*}[t]
\centering
\caption{Evidence hierarchy used throughout the manuscript.}
\label{tab:evidence-tier}
\scriptsize
\begin{tabular}{p{3.2cm}p{3.8cm}p{7.8cm}}
\toprule
Tier & Artifact family & Permitted interpretation \\
\midrule
Confirmatory verification pass & EMBER identity AL/fidelity & Main evidence for RQ2/RQ3; use in abstract and conclusion within dataset scope \\
Original confirmatory protocol & MalBehavD-V1 1,922 & Channel complementarity and component ablation; D does not add stable value \\
Measured development evidence & Fusion-only 1,922 & Main RQ1 comparison; explicitly label same-cohort post-confirmatory development \\
Exploratory development & Robust PID v2, horizon 40 & Hypothesis and algorithm candidate; do not claim external confirmation \\
Exploratory controlled stress & Stateful gated residual PID & Behavior under specified synthetic delay stress only \\
Measured feasibility & Timestamp pilot & Data-source/parser feasibility; no performance or D-term conclusion \\
Excluded & Legacy, simulated, or dashboard values & No use as manuscript evidence \\
\bottomrule
\end{tabular}
\end{table*}

The source-of-truth report identifies experimental commit
\texttt{8bfa06b65b93}. That commit and the new result directories were not present
on the remote branch available to this writing environment. The numerical paper
is therefore assembled from the locked readiness report and its artifact map,
not described as an independent replay. Before submission, the source commit,
artifacts, hash verification, and figure generation must be synchronized in a
clean environment.
